
\subsection{电学实验}

\begin{tabularx}{\columnwidth}{XX}
  电流表改装分流电阻&$R=\dfrac{I_gR_g}{I-I_g}$\\
  电压表改装串联电阻&$R=\dfrac{U}{I_g}-R_g$\\
  欧姆表回路&$E=i(R_{\text{内}}+R_x)$\\
  欧姆表中值电阻&$R_{\text{中}}=R_{\text{内}}$\\
  电桥平衡条件&$\dfrac{R_1}{R_2}=\dfrac{R_3}{R_x}$\\
\end{tabularx}

\begin{formulanotebox}
实验题中公式往往不难，关键是系统误差来自哪里：电流表分压会使测量值偏大，电压表分流会使测量值偏小。

量程选择要兼顾安全和精度，通常让指针偏转达到量程的三分之一以上较合适。
\end{formulanotebox}

\begin{keypointbox}[title={\strut\textcolor{white}{\bfseries 重点知识点：伏安法测电阻}}]
\textbf{电流表内接}：电压表读数包含电流表分压，测得电阻偏大，适合测大电阻。\\

\textbf{电流表外接}：电流表读数包含电压表支路电流，测得电阻偏小，适合测小电阻。\\

\textbf{滑动变阻器}：限流式结构简单、耗能小；分压式可使待测元件两端电压从零开始连续调节。
\end{keypointbox}

\textbf{游标卡尺和螺旋测微器读数}

\begin{itemize}[leftmargin=1.5em]
  \item 游标卡尺读数等于主尺读数加副尺对齐格数乘以精度。
  \item 螺旋测微器读数等于固定刻度读数加可动刻度读数，通常需要估读一位。
  \item 所有读数都要带单位，并注意毫米、厘米之间的换算。
\end{itemize}

\begin{casetypebox}[title={\strut\textcolor{white}{\bfseries 常见题型：实验方案选择}}]
\begin{itemize}[leftmargin=1.5em]
  \item 若要求测伏安特性曲线或电压从零开始调节，优先选分压式。
  \item 若待测电阻较小，电流表外接误差通常较小；若待测电阻较大，电流表内接误差通常较小。
  \item 半偏法测电表内阻时，要判断“半偏前后电源输出是否近似不变”，否则会产生系统误差。
  \item 电桥法利用平衡时检流计示数为零，适合高精度测电阻。
\end{itemize}
\end{casetypebox}
